压缩常驻层次能放入设备、且目标 recall 位于已测竞争区间时，JHQ-GPU 值得评估。
E1 支持 batch 1,000 的高 recall 吞吐；T5 则显示 recall 0.97、batch 8 下，
四个数据集均由 CAGRA fp32 领先。batch 32 的 OpenAI-3072 在插值下接近平手、
实测阈值有利于 JHQ，其余三个有利于 CAGRA。低 recall 优势同样依赖数据集和网格。
分组调度能改善 JHQ 内部执行，但不保证跨系统领先（\cref{sec:7-smallbatch}）。

FULL 是质量优先校准的实测起点，之后应使用部署查询验证并直接计时复用总成本。
应在实际 batch 制度下联合比较质量、总时间和常驻显存。
单客户端同步实验尚未证明并发服务尾延迟、在线更新性能或总体 recall 保证。
